\section{有限域与多项式插值}
\label{sec:04-Discrete-Mathematics/05-Coding-and-Polynomial-Methods/01}

假设你要通过网络发送两个数字——传感器的两个读数。信道偶尔出错，你希望接收方能发现（甚至纠正）错误。最笨的办法是把每个数发三遍，靠投票决定正确值。每传两个有效数就要发六个，冗余两倍。能不能只多传一个数，就获得检错或纠错能力？

换一个来源想同一个问题。平面上给定两个点，恰好有一条直线穿过它们。你告诉接收方两个点的坐标，接收方自己就能算出这条直线在任何位置的取值。现在你再把直线在第三个位置的取值发给对方——接收方核对三个点是否共线，就能发现传输错误。两个点携带信息，第三个点是冗余的检查量。这就是编码的雏形：用代数关系把数据绑定在一起，违背关系的数据就是错的。

但这个想法要真正在计算机里跑起来，我们需要的不只是几何直觉，而是一套精确的代数结构。从几何直觉走到精密编码，中间至少隔着两个坎：数系和次数。

\subsection{为什么不能随便对整数取模}

直觉说``两点确定一条直线''。但在离散世界中——代码的变量只能取有限个值——这个直觉还有效吗？

考虑模 6 的算术：元素是 \(\{0,1,2,3,4,5\}\)，加法和乘法后取模 6。试着解一条直线。假设我们知道 \(f(0)=0\)，\(f(1)=1\)，想求 \(f(2)\)。设 \(f(x)=ax+b\)，由 \(f(0)=0\) 得 \(b=0\)；由 \(f(1)=1\) 得 \(a=1\)。于是 \(f(x)=x\)，\(f(2)=2\)。看起来没问题。

但再加一个点试试。若 \(f(0)=0\)，\(f(1)=0\)，猜直线 \(f(x)=0\)。按两点确定一直线，\(f(2)\) 应该是 \(0\)。好。

现在出难题：\(f(0)=0\)，\(f(1)=0\)，\(f(2)=1\)，而 \(f(x)=ax\)（因为 \(b=0\)）。从 \(f(2)=1\) 得 \(2a\equiv1\pmod6\)。在模 6 下，2 乘以 0,1,2,3,4,5 分别得 0,2,4,0,2,4——永远得不到 1。方程无解。因为 2 和 6 不互素，2 在模 6 下没有乘法逆元，你不能``除以 2''。

更麻烦的事在后面。模 6 时，
\[
2\cdot 3\equiv 0\pmod 6,
\]
而 2 和 3 都不是零。非零元素相乘得零——这样的元素叫零因子。零因子的存在会从根上破坏多项式的基本事实：一个非零多项式可以有很多很多根。例如 \(2x\) 在模 6 下，\(x=0\) 和 \(x=3\) 都让它得零——一次多项式有两个根。

若把模数换成素数 \(p\)，情况彻底不一样了。集合
\[
\mathbb F_p=\{0,1,\ldots,p-1\}
\]
配合模 \(p\) 加法和乘法构成一个域。域的关键性质：每个非零元素 \(a\) 都有唯一乘法逆元 \(a^{-1}\)，满足 \(a\cdot a^{-1}\equiv1\pmod p\)。为什么素数模下必有逆元？因为 \(a<p\) 且 \(p\) 是素数，\(a\) 与 \(p\) 互素；由 Bézout 恒等式，存在 \(x,y\) 使 \(ax+py=1\)，模 \(p\) 后即 \(ax\equiv1\pmod p\)。这个 \(x\) 就是 \(a^{-1}\)（模逆的存在条件与系统求法见\hyperref[sec:04-Discrete-Mathematics/04-Number-Theory/02]{``模运算与模逆''}）。

在 \(\mathbb F_7\) 中，\(3^{-1}=5\)，因为 \(3\cdot5=15\equiv1\pmod7\)。方程 \(3x=4\) 有唯一解 \(x=5\cdot4\equiv6\pmod7\)。除以一个数就是乘以它的逆元——除法回来了，多项式插值需要的全部四则运算在 \(\mathbb F_p\) 中都有了。

动手算几个逆元能建立对这个系统的体感。在 \(\mathbb F_{11}\) 中，要找 \(4^{-1}\)：尝试 \(4\cdot3=12\equiv1\pmod{11}\)，所以 \(4^{-1}=3\)。要找 \(7^{-1}\)：\(7\cdot8=56\equiv1\pmod{11}\)，所以 \(7^{-1}=8\)。对小素数模，逆元可以手工试出来；对大素数模，用扩展 Euclid 算法一步到位。关键不是逆元怎么求，而是每个非零元素都有逆元——这个保证来自素数模，而非一般的合数模。

有限域不只有 \(\mathbb F_p\) 这一种。还存在含 \(p^m\) 个元素的扩域，构造方式涉及多项式商环。但对于理解编码和插值的主线，\(\mathbb F_p\) 作为工作环境已经足够——后面所有论证只需要``域''这个层级的性质。

\subsection{低次多项式为什么有刚性}

在域上，一个非零的 \(d\) 次多项式至多有 \(d\) 个不同的根。证明的逻辑链很短：若 \(f(a)=0\)，用一次式 \(x-a\) 去除 \(f(x)\)：
\[
f(x)=(x-a)q(x)+r,
\]
其中余式 \(r\) 是常数，因为除式是一次式。代入 \(x=a\)：\(f(a)=0+r=r\)，而已知 \(f(a)=0\)，故 \(r=0\)。所以 \(x-a\) 整除 \(f(x)\)，\(f\) 的次数比 \(q\) 的次数多 1。每找到一个不同的根，就剥离一个一次因子，次数减 1。\(d\) 次多项式最多容纳 \(d\) 个互异的一次因子，因此至多有 \(d\) 个不同的根。

``在域上''这个条件就是前面的伏笔。模 8 时（8 不是素数），多项式
\[
f(x)=x^2-1
\]
在 \(x=1,3,5,7\) 处全部得零——一个二次多项式竟有四个不同的根。问题不出在多项式本身，而出在模 8 的算术不是域。2 和 4 都是零因子：\(2\cdot4\equiv0\pmod8\)。没有乘法逆元的数系，连``一次多项式至多有一个根''都保不住。

根的上界直接带来插值唯一性。假设两个次数都小于 \(k\) 的多项式 \(f\) 和 \(g\)，在 \(k\) 个不同位置 \(x_1,\ldots,x_k\) 上取值相同。考虑差
\[
h(x)=f(x)-g(x),
\]
\(h\) 的次数也小于 \(k\)，却拥有 \(k\) 个不同的根 \(x_1,\ldots,x_k\)。除非 \(h\) 是零多项式，否则这与根的上界矛盾。所以 \(h\equiv0\)，即 \(f=g\)。

结论很清楚：
\(k\) 个横坐标互异的点，唯一确定一个次数小于 \(k\) 的多项式。

一条边界说明：这里的约束是``次数小于 \(k\)''，不是``次数不超过 \(k\)''。如果允许次数恰好为 \(k\)，你总可以在任何通过这 \(k\) 个点的多项式上叠加
\[
c\cdot(x-x_1)(x-x_2)\cdots(x-x_k)
\]
的任意倍数——它在那 \(k\) 个点处取零，不改变插值条件，但改变多项式本身。唯一性只在``次数严格小于点数''的范围内成立。自由度一旦正巧等于点数，解不唯一。

\subsection{Lagrange 插值怎样构造这条曲线}

唯一性告诉我们答案只有一个，但没把这个答案交到手里。Lagrange 插值的思路是先造一组``基函数''，每个基函数在自己的插值点取 1，在其余点全取 0：
\[
L_i(x)=\prod_{j\neq i}\frac{x-x_j}{x_i-x_j}.
\]
因为横坐标互异，分母 \(x_i-x_j\) 非零，在域中可以求逆。这个设计很直接：分子在所有 \(x_j\)（\(j\neq i\)）处归零，分母恰好让它在 \(x_i\) 处归一。

有了基函数，给定数据 \((x_i,y_i)\) 时，插值多项式直接写成加权和：
\[
f(x)=\sum_{i=1}^{k}y_i L_i(x).
\]
每个 \(L_i\) 是 \(k-1\) 次的（分子是 \(k-1\) 个一次因子的乘积），求和后次数小于 \(k\)。在 \(x_i\) 处，只有 \(L_i(x_i)=1\)，其余 \(L_j(x_i)=0\)（\(j\neq i\)），所以 \(f(x_i)=y_i\)——恰好满足要求。

来看一个能从头算到尾的例子。在 \(\mathbb F_5\) 中，已知
\[
f(0)=2,\qquad f(1)=4,
\]
并假定 \(f\) 的次数小于 2（即它是一条直线）。两点直接确定直线：
\[
f(x)=2+(4-2)x=2+2x\pmod5.
\]
即使不发送 \(f(2)\)，接收方也知道它应为
\[
f(2)=2+2\cdot2=6\equiv1\pmod5.
\]
前两个取值携带信息，第三个点的取值由直线唯一决定——它没有引入任何新的自由信息。如果收到的三个点是 \((0,2),(1,4),(2,0)\)，接收方算得直线在 \(x=2\) 处应为 1，实际收到 0，三个点不共线，立刻知道传输出了错。

从这个微型例子已经能闻到纠错编码的思路：信息位是两个点的纵坐标，冗余位是第三个点的纵坐标；它们被一个低次多项式的代数关系绑在一起。关系就是``存在一条直线同时穿过三个点''。任何违背这个关系的数据都会被识别为错误。

再来一个有三个点的例子，在 \(\mathbb F_7\) 中。假设要发送三个数 \(1,3,2\)。取插值点 \(x=0,1,2\)，构造次数小于 3 的多项式。先用 Lagrange 基函数：
\[
\begin{aligned}
L_0(x)&=\frac{(x-1)(x-2)}{(0-1)(0-2)}=\frac{(x-1)(x-2)}{2},\\
L_1(x)&=\frac{(x-0)(x-2)}{(1-0)(1-2)}=\frac{x(x-2)}{-1},\\
L_2(x)&=\frac{(x-0)(x-1)}{(2-0)(2-1)}=\frac{x(x-1)}{2}.
\end{aligned}
\]
在 \(\mathbb F_7\) 中，\(2^{-1}=4\)（因为 \(2\cdot4=8\equiv1\)），所以 \(1/2=4\)，\(1/(-1)=-1=6\)。于是
\[
f(x)=1\cdot L_0(x)+3\cdot L_1(x)+2\cdot L_2(x)
= 1 + 3x + 3x^2 \pmod7.
\]
编码时除了原始三个值外，再发送 \(f(3),f(4)\) 等额外校验值。接收方收到任意三个点就能重建多项式，多收到的点用来校验一致性。这正是 Reed--Solomon 码的构造方式——它提供的不只是``存在一条多项式穿过这些点''，而是``如果收到的点多于多项式次数，不一致就是错误''。

\subsection{从插值到秘密共享}

同一个代数结构还能解释另一件看似不相关的事：Shamir 门限秘密共享。设想你要把一个秘密数字分给 5 个人保管，但规定至少 3 人合力才能恢复秘密。

办法：构造一个次数不超过 2 的多项式，把秘密放在常数项 \(f(0)\)，其余两个系数随机选择。然后在某个足够大的有限域中，给第 \(i\) 个人一个点 \((i,f(i))\)。

任意 3 个人凑在一起，握有 3 个不同的点。由插值唯一性，这 3 个点唯一确定那个次数小于 3 的多项式，重建整个 \(f\)，读出秘密 \(f(0)\)。只有 2 个人时，握有 2 个点，而通过这 2 个点的二次多项式并不唯一——二次多项式有 3 个自由系数，2 个点只施加了 2 个约束，还剩 1 维自由空间。在这个自由度里，常数项可以取许多不同的值，秘密没有被确定。

值得把两个机制分开看清楚：
\begin{itemize}
\tightlist
\item
  任意 \(k\) 份能唯一恢复秘密——这是低次多项式插值唯一性的直接推论，是代数刚性：点数等于次数时，多项式被锁死。
\item
  少于 \(k\) 份时秘密不确定——这来自两个来源的叠加。一是点数少于次数时，插值不唯一（自由度大于零）。二是那些自由系数是随机选的——随机性让敌手无法通过枚举来碰运气。
\end{itemize}

用一个微型算例把这两层说清楚。在 \(\mathbb F_{17}\) 中，秘密是常数项 7。随机选一次系数 3，多项式为 \(f(x)=3x+7\)。给三个人的份额是 \(f(1)=10\)，\(f(2)=13\)，\(f(3)=16\)。任意两方——比如拿到 \(f(1)=10\) 和 \(f(2)=13\) 的两方——握有两个点，一次多项式被唯一确定（两点定一直线），恢复秘密 \(f(0)=7\)。只有一方拿到的 \(f(1)=10\)？满足 \(f(1)=10\) 且常数项为任意值的一次多项式有一整个族（斜率可变）——秘密无法确定。

这里恢复靠的是``点数等于次数+1时插值唯一''，隐藏靠的是``点数不足时解空间维度大于零，而随机斜率让枚举无意义''。两个机制的代数根基不同，能力也不同。

把这两个来源混在一起说成``插值保证了安全性''是不准确的。插值唯一性保证恢复，插值不唯一性加上随机系数保证信息隐藏。两个机制各司其职，各自依赖不同的代数事实。

\subsection{插值的另一种视角：Vandermonde 方程组}

把插值写成线性方程组，能更清楚地看见``域''和``次数''各自扮演什么角色。给定 \(k\) 个点 \((x_i,y_i)\)，要找次数小于 \(k\) 的多项式 \(f(x)=c_0+c_1x+\cdots+c_{k-1}x^{k-1}\)，使 \(f(x_i)=y_i\)。展开成矩阵形式：
\[
\begin{pmatrix}
1 & x_1 & x_1^2 & \cdots & x_1^{k-1}\\
1 & x_2 & x_2^2 & \cdots & x_2^{k-1}\\
\vdots & \vdots & \vdots & \ddots & \vdots\\
1 & x_k & x_k^2 & \cdots & x_k^{k-1}
\end{pmatrix}
\begin{pmatrix}c_0\\c_1\\\vdots\\c_{k-1}\end{pmatrix}
=
\begin{pmatrix}y_1\\y_2\\\vdots\\y_k\end{pmatrix}.
\]
这个系数矩阵就是 Vandermonde 矩阵。当 \(x_i\) 互异时，它的行列式非零（在域中），所以系数向量唯一存在。Vandermonde 行列式非零的条件是``所有 \(x_i-x_j\neq0\)''——这恰好要求每个非零差都有逆元，即工作环境必须是域。

Lagrange 插值和 Vandermonde 方程组是同一枚硬币的两面。Lagrange 直接构造出多项式（避免了求逆矩阵），Vandermonde 视角则暴露了背后的线性代数本质：插值就是在一个基底下解线性方程组。当工作环境不是域时，方程组可能无解或有无穷多解，并不是 Lagrange 公式本身失效——而是除法根本做不了。

\subsection{回头想想：实数插值和有限域插值差在哪里}

在实数上做插值——中小学的 Lagrange 插值——数值可以连续变化，曲线可以画在坐标纸上。在有限域上做插值，数值只在 \(\{0,1,\ldots,p-1\}\) 里跳，没有``连续''这个概念的容身之处。

但代数骨架一模一样。Vandermonde 矩阵、Lagrange 基函数、根的上界、插值唯一性——这些在实数和有限域中的表述和证明完全相同，区别只在于逆元从哪里来。实数中的逆元是 \(1/(x_i-x_j)\)，有限域中同样还是 \(1/(x_i-x_j)\)——只不过这个倒数现在指的是模 \(p\) 下的乘法逆元。

实数插值追求的是光滑性、逼近精度、Runge 现象。有限域插值完全不关心中间值长什么样——它只关心在离散的评估点上等式是否成立。实数插值的误差是连续的逼近误差，有限域插值的``误差''是离散的 Hamming 距离——根本不在同一个度量空间里。

把这两个语境下的插值放在一起看，你会发现代数结构（域公理 + 多项式环）比具体数值更重要。一旦代数结构定下来了，插值的全部性质随之锁定。具体是实数域还是有限域，只影响你怎么可视化结果——不影响定理的内容。

\subsection{这一节留下了什么}

有限域给了我们一个能做除法的离散数系，低次多项式则把少数自由参数扩展成一串相互制约的取值——这就是编码的代数基础设施。从几何直觉（两点定直线）出发，经过模运算的陷阱和域的筛选，落脚在多项式的根上界这个精确事实上。再把插值翻译成 Vandermonde 方程组，线性代数给了我们另一个确认：只要横坐标互异且系数来自域，解就唯一存在。

但这一节还没有回答真正的问题：一组合法的编码字符串彼此需要相隔多远，才能保证发现或纠正一定数量的传输错误？这是下一篇要处理的事，它的答案——Hamming 距离和 Singleton 界——与这里多项式的次数上界有着紧密的代数联系。多项式插值唯一性，正是 Singleton 界在 Reed--Solomon 码这个特殊构造中的体现。从这里出发，两件东西将被连在一起：多项式的次数（决定了有多少个``自由系数''）和编码的最小距离（决定了能纠正多少错误）。
